\section{Claim Matrix}
\label{sec:claims}

The paper's claims are intentionally narrower than the simulator's design
space. Table~\ref{tab:claim-matrix} separates claims already supported by the
current empirical package from hypotheses that are unsupported or still open.
This keeps the contribution focused on the observed information-safety
tradeoff rather than on a generic claim that communication always improves
evacuation.

\begin{table*}[!t]
\centering
\caption{Claim matrix for the current Chiyoda paper. The primary evidence is
the 50-seed information-control study, with the 30-seed ablation and
message-quality studies and the 900-run regime robustness grid used as
supporting stress tests.}
\label{tab:claim-matrix}
\begin{tabular}{p{0.18\textwidth}p{0.18\textwidth}p{0.34\textwidth}p{0.20\textwidth}}
\toprule
Claim & Status & Evidence & Paper role \\
\midrule
Emergency communication is a safety-control action, not merely an
entropy-reduction input. & Supported & Active policies change belief accuracy,
entropy, exposure pressure, queue pressure, and harmful convergence in
different directions. & Central thesis and framing. \\
\addlinespace
Static/local messaging has the best information-safety efficiency in the
main study. & Supported & Static beacons rank first by ISE in the 50-seed
main comparison, ahead of global broadcast and all adaptive policies. & Main
positive result. \\
\addlinespace
Broad reach is not equivalent to useful safety effect. & Supported & Global
broadcast reaches many agents and slightly improves some physical aggregates,
but ranks below local and capacity-aware policies by ISE. & Distinguishes the
paper from broadcast-more baselines. \\
\addlinespace
Entropy or accuracy gains can create harmful convergence. & Supported & The
main and ablation studies show that stronger belief improvement can coexist
with queue pressure, exit concentration, and weak evacuation-count gains. &
Mechanistic contribution and reviewer-facing caution. \\
\addlinespace
Exposure-aware targeting clearly reduces hazard exposure. & Not supported &
Exposure-aware targeting is close to baseline exposure and is not the lowest
exposure policy in the 50-seed main comparison. & Negative result; avoid
overclaiming. \\
\addlinespace
Static/local messaging is robust by information-safety efficiency across
hazard severity and familiarity regimes. & Supported & In the 900-run regime
grid, static beacons have the highest ISE in every hazard--familiarity cell. &
External-validity support for the main positive result. \\
\addlinespace
The policy that maximizes completed evacuations is stable across regimes. &
Not supported & The best evacuation-count policy changes across the
robustness grid, including static beacons, global broadcast, entropy targeting,
and no intervention. & Separates safety-efficiency claims from operational
winner claims. \\
\addlinespace
Validated LLM guidance is best understood as sparse safety control, not as a
replacement for deterministic communication. & Supported as bounded extension &
Live OpenAI prompt, regime, target-selection, and budget-equivalence studies
show high sparse ISE, prompt-sensitive HCI tradeoffs, exact replay, and weaker
ISE under equal generated-message budgets. & Larger extension result; does not
supersede deterministic baseline claims. \\
\bottomrule
\end{tabular}
\end{table*}

The final rows are deliberately separate. The robustness grid strengthens the
static-beacon result for information-safety efficiency, but it does not make a
universal operational prescription. The LLM extension is also bounded: it
shows that generated guidance can be evaluated as safety control, while ruling
out the stronger claim that richer generated language automatically improves
harmful convergence or replaces conservative deterministic communication.
Completed evacuations remain noisy and regime dependent, which is consistent
with the paper's central claim: information control must be evaluated by
coupled belief and safety effects rather than by a single aggregate endpoint.
